\chapter{Fazit}

EU hat das Thema Digitalisierung als politische Priorität. Es ist eine wichtige Voraussetzung für die wirtschaftliche und gesellschaftliche Entwicklung der EU. Dabei ist Cloud als zentraler Baustein der Digitalisierung fungiert. Ohne die Skalierbarkeit, Flexibilität und Kosteneffizienz von Cloud-Diensten wäre die digitale Transformation schwierig vorstellbar, besonders für kleinere Unternehmen.

Allerdings bringt Cloud Computing einen fundamentalen Widersprucht mit sich: Es setzt voraus, dass Kontrolle über digitale Infrastruktur an externe Anbieter abgegeben wird, während digitale Souveränität diese Kontrolle zurückfordert. In der EU gibt es einen starken Abhängigkeiten von US-amerikanischen Cloud-Diensten. Rund 70 Prozent der Marktanteile werden nur von 3 US-Hyperscaler dominiert, die einer fremden Rechtsordnung unterliegen. Der Cloud-Act ermöglicht US-Regierung den Zugriff auf Daten , unabhängig davon, wo sie gespeichert sind. Technischer Vendor Lock-In erschwert den Wechsel zu alternativen Anbietern. Die letzte geopolitische Entwicklungen hat gezeigt, dass die Abhängigkeiten als geopolitische Druckmittel gegen die EU selbst eingesetzt werden.

Die bestehende Ansätze der EU sowohl regulatorisch (DSGVO, DMA, DSA, Data Act), als auch technisch (Gaia-X, Cloud Sovereignty Framework) adressieren diese Probleme teilweise, aber nicht vollständig. Die Asymmetrie der Marktmacht, die technische Herausforderung zum Wechsel der Anbietern und die rechtliche Unsicherheit zum Jurisdiktionsproblem bleiben bestehen.
Zudem sind die souveränen Angeboten der US-Hyperscaler immer noch unklar und geben keine Garantie, dass sie tatsächlich die Kontrolle über die digitale Infrastruktur an der EU zurückgeben. 

Komplette Souveränität im Cloud-Bereich ist theoretisch möglich, China hat dies durch gesetzlich erzwungene Eigentumsübertragung an lokale Betreiber bewiesen. Allerdings wäre dieser Ansatz für die EU nicht praktikabel, da erhebliche Investitionen für die einheimische digitale Infrastuktur und menschliche Ressourcen erforderlich wären. Ein umsetzbarerer Ansatz wäre die \textit{selektive Souveränität}: Die EU bevorzugen, kritische Infrastrukturen ausschließlich bei europäischen Cloud-Anbietern oder lokal im eigenen Rechenzentren (Hybrid-Cloud-Ansatz mit Open-Source-Technologien) zu betrieben und client-seitige Verschlüsselung für besonders sensitive Daten z.B persönliche Daten verpflichtend einzuführen. Die Verschlüsselung kann zwar vor dem Zugriff durch US-Behörden schützen, aber sie beseitigt die Abhängigkeit von US-Anbietern nicht, damit werden sensitive Daten trotz Einsatz von US-Cloud-Diensten geschützt.

Das Ziel digitaler Souveränität sollte nicht vollständige technologische Unabhängigkeit sein. Dies wäre weder realistisch noch 
wirtschaftlich sinnvoll. Das Ziel muss sein, dass die EU nie so abhängig wird, dass sie erpressbar ist. Souveränität bedeutet nicht Autarkie, sondern glaubwürdige Alternativen zu haben und die eigene Abhängigkeit bewusst und möglich zu gestalten. 